\documentclass[10pt]{article}
\usepackage[margin=0.8in]{geometry}
\usepackage{xcolor}
\usepackage[colorlinks=true,allcolors=blue!55!black]{hyperref}
\pagestyle{empty}
\setlength{\parindent}{0pt}
\setlength{\parskip}{0.7em}

\begin{document}

\begin{flushright}
Signal and Image Processing Institute\\
Ming Hsieh Department of Electrical and Computer Engineering\\
University of Southern California\\
Los Angeles, California, USA\\
\texttt{ajoshi@usc.edu}
\end{flushright}

Editors\\
\emph{Imaging Neuroscience}

Dear Editors,

Please consider the Research Article ``Contralateral-Normalized Deformation and
Stationary Log-Velocity Decomposition after Chronic Stroke: Geometric
Characterization and Aphasia Prediction in the Aphasia Recovery Cohort'' for
publication in \emph{Imaging Neuroscience}.

This study addresses a recurring interpretive problem in lesion neuroimaging:
registration fields near chronic lesions can be scientifically useful without
being physical measurements of tissue displacement or pressure. We introduce a
contralateral-normalized, lesional-only deformation proxy, construct a
positive-Jacobian stationary log-velocity embedding, and apply an explicitly
bounded Helmholtz--Hodge decomposition. Every log-domain case is audited by
exponentiating the estimated velocity and comparing the reconstruction with its
input deformation.

The main scientific result is deliberately negative and useful. In 210
participants from the openly available Aphasia Recovery Cohort, deformation
features produced only a small and uncertain reduction in held-out Western
Aphasia Battery Aphasia Quotient error. Hodge descriptors were descriptively
associated with aphasia severity but did not improve nested cross-validated
prediction beyond conventional lesion or displacement features. The manuscript
distinguishes this geometric analysis from physical pressure inference and
reports the numerical, registration-sensitivity, and statistical checks needed
to support that boundary.

The work fits \emph{Imaging Neuroscience} because it combines open human
neuroimaging, deformable registration, reproducible computational methodology,
and out-of-sample brain--behavior modeling. Code, tests, aggregate results, and
LaTeX sources are publicly available in the
\href{https://github.com/ajoshiusc/ARCDeformationAnalysis}{ARCDeformationAnalysis
repository}; the source dataset is available from OpenNeuro.

OpenAI Codex assisted with code refactoring, statistical implementation,
reproducibility checks, and language editing. The author reviewed the analysis
and accepts responsibility for the submitted code and text. The corresponding
disclosure appears in the manuscript.

Before submission, the author will confirm originality and exclusive
consideration and supply the final funding and competing-interest declarations
in the manuscript and submission system.

Sincerely,

\vspace{0.5em}
Anand A. Joshi\\
Signal and Image Processing Institute\\
University of Southern California\\
\texttt{ajoshi@usc.edu}

\end{document}
